\subsubsection*{Orientation}

\begin{tcolorbox}[
  colback=gray!6,
  colframe=gray!40,
  arc=2pt,
  left=8pt, right=8pt, top=6pt, bottom=6pt,
  title={\small\textbf{Toolkit}},
  fonttitle=\small\bfseries
]
This orientation uses logical equivalence, truth assignments, literals, clauses,
and algebraic rewrite laws to motivate normal forms.
\end{tcolorbox}

\begin{remark*}[Exposition]
A normal form is a controlled syntactic shape for a propositional formula.
Instead of allowing an arbitrary mixture of connectives, parentheses, and
nested subformulas, a normal-form convention specifies which connectives may
appear and how they may be arranged. The point is not to change the proposition
being expressed, but to put the proposition into a form whose structure is easy
to inspect.

Normal-form conversion is therefore governed by logical equivalence. When a
formula is converted into a normal form, the converted formula must have the
same truth value as the original formula under every truth assignment. The
algebra of propositions supplies the equivalence laws used to justify these
conversions, while the semantic theory explains why preserving truth under
every assignment is the right standard.
\end{remark*}

\begin{remark*}[Section map]
This section uses the preceding syntax, semantics, and algebra of propositions
as its starting point. It first isolates literals and clauses, then introduces
negation normal form, disjunctive normal form, conjunctive normal form,
canonical normal forms, and semantic classification uses. Later arguments
will use these forms to organize satisfiability questions, proof search,
functional completeness, and resolution-style methods.
\end{remark*}

\begin{remark*}[Why normal forms matter]
Normal forms trade arbitrary expression for regular structure. In a truth-table
argument, regular structure helps classify formulas as tautologies,
contradictions, satisfiable formulas, or contingencies. In proof search and
satisfiability procedures, regular structure turns a formula into a finite
combinatorial object that can be searched. In functional completeness, normal
forms show how complicated truth functions can be built from simpler
connectives.
\end{remark*}
